% GENERATED FILE. Edit canonical lessons, then run npm run generate:books.
% canonical-source-hash: fnv1a64:b44d0a2ddec1148b
% canonical-lessons: MR-C440-duskal, MR-C440-dhagphuti, MR-C440-kadaka, MR-C440-jhodapne, MR-C440-osarne

\chapter{Droughts and Cloudbursts}
\label{ch:mr-droughts-and-cloudbursts}
\hlchaptermodality{440}

\begin{chapteropening}
\textbf{By the end of this chapter:} \emph{I can talk about droughts, cloudbursts and hard spells, and say rain lashed and the water went down.}

Complete the last lesson of chapter 440: I can talk about droughts, cloudbursts and hard spells, and say rain lashed and the water went down.
\end{chapteropening}

\section[\texorpdfstring{\mr{दुष्काळ} (duṣkāḷ)}{दुष्काळ (duṣkāḷ)}]{\mr{दुष्काळ} (duṣkāḷ) — a drought; a famine}

\label{lesson:MR-C440-duskal}



\begin{quote}
Before the new one: say the Marathi for to howl, then the Marathi for to crackle.
\end{quote}

\subsection*{You'll want to know: \mr{दुष्काळ}}
\textbf{\mr{दुष्काळ}} — \emph{duṣkāḷ} — \textquotedblleft{}a drought; a famine\textquotedblright{}.

\begin{quote}
\mr{त्या} \mr{वर्षी} \mr{मोठा} \mr{दुष्काळ} \mr{पडला}.
\end{quote}

\begin{quote}
\emph{tyā varṣī moṭhā duṣkāḷ paḍlā.}
\end{quote}

\begin{quote}
That year there was a great drought.
\end{quote}

Masculine.

\mr{दुष्काळ} \mr{पडणे}, a drought strikes.

\begin{cousinweb}
\textbf{\mr{दुष्काळ}} (\emph{duṣkāḷ}) is a tatsama, taken from Sanskrit \emph{duṣkāla}, \textquotedblleft{}a bad time, a famine\textquotedblright{} (CDIAL 6470), where CDIAL also lists the inherited Marathi \emph{dukāḷ}.
\end{cousinweb}

\begin{grammarlens}[title={What you've built}]
One more word for getting around.
\end{grammarlens}

\subsection*{Your turn}
\begin{itemize}
\raggedright
  \item \emph{Say it:} \emph{duṣkāḷ}
  \item \emph{Say it:} \emph{duṣkāḷ}, once more
  \item \emph{Say it:} \emph{tyā varṣī moṭhā duṣkāḷ paḍlā.} — That year there was a great drought.
  \item \emph{Recall:} say the Marathi for to howl, then the Marathi for to crackle, then say \emph{duṣkāḷ} again
\end{itemize}

\begin{tcolorbox}[breakable,colback=teal!4,colframe=teal!35!black,title={Before you move on}]
What does \mr{दुष्काळ} mean? (\textquotedblleft{}A drought; a famine\textquotedblright{}.) Say it once more.
\end{tcolorbox}
\section[\texorpdfstring{\mr{ढगफुटी} (ḍhagphuṭī)}{ढगफुटी (ḍhagphuṭī)}]{\mr{ढगफुटी} (ḍhagphuṭī) — a cloudburst}

\label{lesson:MR-C440-dhagphuti}



\begin{quote}
Before the new one: say the Marathi for to crackle, then the Marathi for a drought.
\end{quote}

\subsection*{You'll want to know: \mr{ढगफुटी}}
\textbf{\mr{ढगफुटी}} — \emph{ḍhagphuṭī} — \textquotedblleft{}a cloudburst\textquotedblright{}.

\begin{quote}
\mr{डोंगरावर} \mr{ढगफुटी} \mr{झाली}.
\end{quote}

\begin{quote}
\emph{ḍoṅgarāvar ḍhagphuṭī jhālī.}
\end{quote}

\begin{quote}
There was a cloudburst on the mountain.
\end{quote}

Feminine.

\mr{ढग}, a cloud, and \mr{फुटणे}, to burst.

\begin{cousinweb}
\textbf{\mr{ढगफुटी}} (\emph{ḍhagphuṭī}) is built from \textbf{\mr{ढग}}, \textquotedblleft{}a cloud\textquotedblright{}, and the stem of \textbf{\mr{फुटणे}}, \textquotedblleft{}to burst\textquotedblright{} (CDIAL 13845).
\end{cousinweb}

\begin{grammarlens}[title={What you've built}]
One more word for getting around.
\end{grammarlens}

\subsection*{Your turn}
\begin{itemize}
\raggedright
  \item \emph{Say it:} \emph{ḍhagphuṭī}
  \item \emph{Say it:} \emph{ḍhagphuṭī}, once more
  \item \emph{Say it:} \emph{ḍoṅgarāvar ḍhagphuṭī jhālī.} — There was a cloudburst on the mountain.
  \item \emph{Recall:} say the Marathi for to crackle, then the Marathi for a drought, then say \emph{ḍhagphuṭī} again
\end{itemize}

\begin{tcolorbox}[breakable,colback=teal!4,colframe=teal!35!black,title={Before you move on}]
What does \mr{ढगफुटी} mean? (\textquotedblleft{}A cloudburst\textquotedblright{}.) Say it once more.
\end{tcolorbox}
\section[\texorpdfstring{\mr{कडाका} (kaḍākā)}{कडाका (kaḍākā)}]{\mr{कडाका} (kaḍākā) — a severe spell}

\label{lesson:MR-C440-kadaka}



\begin{quote}
Before the new one: say the Marathi for a drought, then the Marathi for a cloudburst.
\end{quote}

\subsection*{You'll want to know: \mr{कडाका}}
\textbf{\mr{कडाका}} — \emph{kaḍākā} — \textquotedblleft{}a severe spell\textquotedblright{}.

\begin{quote}
\mr{यंदा} \mr{थंडीचा} \mr{कडाका} \mr{जास्त} \mr{आहे}.
\end{quote}

\begin{quote}
\emph{yandā thaṇḍīcā kaḍākā jāsta āhe.}
\end{quote}

\begin{quote}
This year the cold is especially severe.
\end{quote}

Masculine.

Of cold or heat: \mr{थंडीचा} \mr{कडाका}, \mr{उन्हाचा} \mr{कडाका}.

\begin{cousinweb}
\textbf{\mr{कडाका}} (\emph{kaḍākā}) is inherited, from the family of Sanskrit \emph{kaṭakaṭā}, where CDIAL lists Marathi \emph{kaḍakṇẽ}, \textquotedblleft{}to crackle\textquotedblright{} (CDIAL 2632): cold sharp enough to crack.
\end{cousinweb}

\begin{grammarlens}[title={What you've built}]
One more word for getting around.
\end{grammarlens}

\subsection*{Your turn}
\begin{itemize}
\raggedright
  \item \emph{Say it:} \emph{kaḍākā}
  \item \emph{Say it:} \emph{kaḍākā}, once more
  \item \emph{Say it:} \emph{yandā thaṇḍīcā kaḍākā jāsta āhe.} — This year the cold is especially severe.
  \item \emph{Recall:} say the Marathi for a drought, then the Marathi for a cloudburst, then say \emph{kaḍākā} again
\end{itemize}

\begin{tcolorbox}[breakable,colback=teal!4,colframe=teal!35!black,title={Before you move on}]
What does \mr{कडाका} mean? (\textquotedblleft{}A severe spell\textquotedblright{}.) Say it once more.
\end{tcolorbox}
\section[\texorpdfstring{\mr{झोडपणे} (jhoḍapṇe)}{झोडपणे (jhoḍapṇe)}]{\mr{झोडपणे} (jhoḍapṇe) — to lash; to batter}

\label{lesson:MR-C440-jhodapne}



\begin{quote}
Before the new one: say the Marathi for a cloudburst, then the Marathi for a severe spell.
\end{quote}

\subsection*{You'll want to know: \mr{झोडपणे}}
\textbf{\mr{झोडपणे}} — \emph{jhoḍapṇe} — \textquotedblleft{}to lash; to batter\textquotedblright{}.

\begin{quote}
\mr{पावसाने} \mr{गावाला} \mr{झोडपलं}.
\end{quote}

\begin{quote}
\emph{pāvsāne gāvālā jhoḍaplaṁ.}
\end{quote}

\begin{quote}
The rain lashed the village.
\end{quote}

Past, with -\mr{ने} (he …): \mr{त्याने} \mr{झोडपलं}, \emph{tyāne jhoḍaplaṁ}.

The rain is the agent and takes -\mr{ने}: \mr{पावसाने} \mr{झोडपलं}.

\begin{cousinweb}
\textbf{\mr{झोडपणे}} (\emph{jhoḍapṇe}) is a Marathi verb for heavy rain or blows beating down; its origin is not settled.
\end{cousinweb}

\begin{grammarlens}[title={What you've built}]
One more word for getting around.
\end{grammarlens}

\subsection*{Your turn}
\begin{itemize}
\raggedright
  \item \emph{Say it:} \emph{jhoḍapṇe}
  \item \emph{Say it:} \emph{jhoḍapṇe}, once more
  \item \emph{Say it:} \emph{pāvsāne gāvālā jhoḍaplaṁ.} — The rain lashed the village.
  \item \emph{Recall:} say the Marathi for a cloudburst, then the Marathi for a severe spell, then say \emph{jhoḍapṇe} again
\end{itemize}

\begin{tcolorbox}[breakable,colback=teal!4,colframe=teal!35!black,title={Before you move on}]
What does \mr{झोडपणे} mean? (\textquotedblleft{}To lash; to batter\textquotedblright{}.) Say it once more.
\end{tcolorbox}
\section[\texorpdfstring{\mr{ओसरणे} (osarṇe)}{ओसरणे (osarṇe)}]{\mr{ओसरणे} (osarṇe) — to recede; to subside}

\label{lesson:MR-C440-osarne}



\begin{quote}
Before the new one: say the Marathi for a severe spell, then the Marathi for to lash.
\end{quote}

\subsection*{You'll want to know: \mr{ओसरणे}}
\textbf{\mr{ओसरणे}} — \emph{osarṇe} — \textquotedblleft{}to recede; to subside\textquotedblright{}.

\begin{quote}
\mr{दोन} \mr{दिवसांनी} \mr{पूर} \mr{ओसरला}.
\end{quote}

\begin{quote}
\emph{don divasānnī pūr osarlā.}
\end{quote}

\begin{quote}
After two days the flood went down.
\end{quote}

Past (he …): \mr{तो} \mr{ओसरला}, \emph{to osarlā}.

Floods, crowds and fevers \mr{ओसरतात}; not the same word as \mr{ओसरी}, a veranda.

\begin{cousinweb}
\textbf{\mr{ओसरणे}} (\emph{osarṇe}) is inherited: CDIAL lists Marathi \emph{osarṇẽ}, \textquotedblleft{}to flow away, decrease, subside\textquotedblright{}, under Sanskrit \emph{apasarati}, \textquotedblleft{}slips off, goes away\textquotedblright{} (CDIAL 465).
\end{cousinweb}

\begin{grammarlens}[title={What you've built}]
One more word for getting around.
\end{grammarlens}

\subsection*{Your turn}
\begin{itemize}
\raggedright
  \item \emph{Say it:} \emph{osarṇe}
  \item \emph{Say it:} \emph{osarṇe}, once more
  \item \emph{Say it:} \emph{don divasānnī pūr osarlā.} — After two days the flood went down.
  \item \emph{Recall:} say the Marathi for a severe spell, then the Marathi for to lash, then say \emph{osarṇe} again
\end{itemize}

\begin{tcolorbox}[breakable,colback=teal!4,colframe=teal!35!black,title={Before you move on}]
What does \mr{ओसरणे} mean? (\textquotedblleft{}To recede; to subside\textquotedblright{}.) Say it once more.
\end{tcolorbox}
